\section{Simultaneous contraction and escape}
\label{sec:simultaneous-contraction}

Proposition~\ref{prop:trapped-row-space} identifies the syndromes whose
entire terminal row space is trapped.  We now treat every other syndrome
without choosing a sequence of intermediate good fibers.  All contracted roots are selected at
once, and only the terminal R5 pencil is counted.

The contraction identity and finite-field grid lemma below are elementary.
The contribution is their coupling to the recursive carrier through one
terminal selector: the selector replaces all intermediate-redundancy casework,
while the final enumerative input is the redundancy-five split-member count
and branch budget of
Section~\ref{sec:r5}.  Its relation to the twisted-cubic line literature is
spelled out in Remark~\ref{rem:r5-line-literature}.

Let \(E\) be two-dimensional over a field \(F\).  For \(0\leq m\leq r-2\),
\(f\in\Gamma^{r-1}E\), and \(R\in\Sym^m(E^\vee)\), define
\(\iota_Rf\in\Gamma^{r-1-m}E\) by
\begin{equation}\label{eq:composite-contraction}
 \langle\iota_Rf,h\rangle=\langle f,Rh\rangle
 \qquad\bigl(h\in\Sym^{r-1-m}E^\vee\bigr).
\end{equation}

\begin{proposition}[Composite contraction]
\label{thm:composite-contraction}%
\coverage{absent}%
For every \(h\in\Sym^{r-2-m}E^\vee\),
\[
 h\in W_{\iota_Rf}\quad\Longleftrightarrow\quad Rh\in W_f.
\]
The construction commutes with base change and the natural
\(\operatorname{GL}(E)\)-action.  If \(R\) is split and squarefree, then
\(Rh\) is split and squarefree exactly when \(h\) is split and squarefree and
has no root in common with \(R\).
\end{proposition}

\begin{proof}
For every \(\lambda\in E^\vee\),
\[
 \langle\iota_Rf,\lambda h\rangle
   =\langle f,R\lambda h\rangle.
\]
These are precisely the two Hankel equations defining the two witness
systems.  The pairing also proves base-change and equivariance.  The final
assertion is the usual coprimality criterion after passage to an algebraic
closure.
\end{proof}

We use one elementary finite-field lemma to choose the roots of \(R\).

\begin{lemma}[Vandermonde grid]
\label{lem:vandermonde-grid}%
\coverage{absent}%
Let \(P\in\F_q[x_1,\ldots,x_m]\) be nonzero and satisfy
\(\deg_{x_i}P\leq d\) for every \(i\).  If \(q>d+m-1+s\), then there are
pairwise distinct \(a_1,\ldots,a_m\in\F_q\), outside any prescribed set of
\(s\) field elements, for which \(P(a_1,\ldots,a_m)\ne0\).
\end{lemma}

\begin{proof}
Multiply \(P\) by
\[
 \prod_{i<j}(x_i-x_j)
 \prod_{i=1}^m\prod_{a\in A_0}(x_i-a),
\]
where \(A_0\subset\F_q\) is the forbidden set.  The product is nonzero and
has degree at most \(d+m-1+s<q\) in each variable.  A polynomial reduced to
degree below \(q\) in every variable cannot vanish on all of \(\F_q^m\).
At a nonzero value, the coordinates are distinct and avoid \(A_0\).
\end{proof}

Put \(m=r-5\).  If \(R=\lambda_1\cdots\lambda_m\), the coefficient row of
\(R\) sends \(f\) to the bottom syndrome
\[
 \iota_Rf=h\operatorname{Cat}_{m,4}(f).
\]

\begin{lemma}[Terminal selector]
\label{lem:terminal-selector}%
\coverage{absent}%
Let \(E\) be defined over \(\F_q\), let
\(0\ne f\in\Gamma^{r-1}E\), and let \(L_f\) be the projectivized row space
of \(\operatorname{Cat}_{m,4}(f)\).  If
\([f]\notin\cE_{r,p}\), there is a homogeneous polynomial
\(F_f\) over \(\F_q\), of degree at most six (at most four when \(p=2\)),
which vanishes on the reduced terminal carrier but whose restriction to
\(L_f\) is not the zero polynomial.
Consequently
\[
 S_f(\lambda_1,\ldots,\lambda_m)
 =F_f(\iota_{\lambda_1\cdots\lambda_m}f)
\]
is a nonzero multihomogeneous \(\F_q\)-polynomial.  On an affine marker chart,
its dehomogenization has degree at most six, respectively four, in each marker
coordinate.
\end{lemma}

\begin{proof}
Proposition~\ref{prop:trapped-row-space} says that \(L_f\) is
not contained in the reduced terminal carrier.  By
Proposition~\ref{prop:reduced-terminal-carrier}, that carrier is
\(V(D)\cup V(I_{\mathrm{res},p})\), where \(D\) has degree three and
\(I_{\mathrm{res},p}\) has generators over the prime field of degree at
most three, or degree one when \(p=2\).  The irreducible space \(L_f\) is
contained in neither component.  Thus \(D|_{L_f}\ne0\), and some listed
generator \(G\in I_{\mathrm{res},p}\) has \(G|_{L_f}\ne0\).  Since the
coordinate ring of \(L_f\) is a domain, \(F_f=DG\) is nonzero on \(L_f\),
vanishes on both components, and has the asserted degree.  It is defined
over \(\F_q\).

Over an algebraic closure, the ordered root-product map on affine cones is
surjective onto \(\Sym^mE^\vee\), and the linear catalecticant map is
surjective onto the affine cone over \(L_f\).  Hence the pullback of the
nonzero restriction \(F_f|_{L_f}\) is nonzero.  Every
coefficient of \(\lambda_1\cdots\lambda_m\) is linear in each marker;
therefore composition preserves the degree bound separately in every
marker.
\end{proof}

\begin{theorem}[Simultaneous-marker escape]
\label{thm:simultaneous-marker-escape}%
\coverage{absent}%
\uses{prop:trapped-row-space,prop:reduced-terminal-carrier,prop:r5-count,lem:r5-branch}%
Let \(E\) be defined over \(\F_q\), let \(r\geq6\), let
\(A\subseteq\PP^1(\F_q)\) have size \(s\), and put
\(S=\PP^1(\F_q)\setminus A\) and
\[
 Q^*_{r,s}=6(r+s)-16+
 \left\lfloor2\sqrt{6(r+s)-18}\right\rfloor .
\]
If \(q\geq Q^*_{r,s}\), \(0\ne f\in\Gamma^{r-1}E\), and
\([f]\notin\cE_{r,p}\), where \(\cE_{r,p}\) is the locus
\eqref{eq:exceptional-locus}, then \(W_f\) contains a split
squarefree form of degree \(r-2\) supported on \(S\).  In characteristic
two it is enough that
\[
 q\geq6(r+s)-22+
 \left\lfloor2\sqrt{6(r+s)-24}\right\rfloor .
\]
\end{theorem}

\begin{proof}
Choose a retained point of \(S\) as infinity by a projective coordinate
change.  Then every deleted point lies in the affine marker chart.
Lemma~\ref{lem:terminal-selector} and
Lemma~\ref{lem:vandermonde-grid} give a split squarefree marker form \(R\),
disjoint from \(A\), for which \(b=\iota_Rf\) lies outside both components of
the terminal carrier.  The required selector inequality is
\(q>r+s\), or \(q>r+s-2\) in characteristic two; it is weaker than the
displayed hypotheses.

Because \(b\notin V(D)=\cP_5\),
Lemma~\ref{lem:basepointfree-off-persistent} makes \(W_b\) base-point-free.
Proposition~\ref{prop:reduced-terminal-carrier} then places its cubic map in
the separable \(S_3\) case.  Proposition~\ref{prop:r5-count} and its branch
budget give at least
\[
 \frac{q+1-2\sqrt q-B_p}{6},\qquad
 B_p=\begin{cases}6,&p=2,\\12,&p\ne2,
 \end{cases}
\]
split squarefree pencil members.  A prescribed rational root occurs in at
most one member of a base-point-free pencil.  Hence a split cubic avoids the
\(m+s\) roots of \(R\) and \(A\) whenever
\[
 q+1-2\sqrt q>B_p+6(m+s).
\]
For an integer \(\Delta\geq0\), the least integer bound forced by
\(q+1-2\sqrt q>\Delta\) is
\[
 \Delta+2+\lfloor2\sqrt\Delta\rfloor.
\]
Taking \(\Delta=6(r+s)-18\), or \(6(r+s)-24\) in characteristic two,
gives the stated thresholds.  Proposition~\ref{thm:composite-contraction}
multiplies the terminal cubic by \(R\), producing the required locator in
\(W_f\).
\end{proof}

The proof yields many locators rather than one isolated witness.

\begin{corollary}[Split-witness abundance]
\label{cor:split-witness-abundance}%
\coverage{absent}%
Fix \(r,s\).  As \(q\) tends to infinity through prime powers satisfying the
hypotheses of
Theorem~\ref{thm:simultaneous-marker-escape}, every syndrome outside
\(\cE_{r,p}\) has at least
\[
 \frac{q^{r-4}}{(r-2)!}-O_{r,s}(q^{r-9/2})
\]
distinct projective split squarefree members of \(W_f\) supported on \(S\),
equivalently distinct unordered root sets.  This denotes a lower bound, not an
asymptotic equality.
\end{corollary}

\begin{proof}
On the affine marker chart, the selector has total degree at most \(6m\).
After deleting diagonals and forbidden values, the number of good ordered
marker tuples is at least
\[
 (q-s)_m-6m q^{m-1},
\]
where \((q-s)_m=(q-s)(q-s-1)\cdots(q-s-m+1)\),
with \(6m\) replaced by \(4m\) in characteristic two.  Each has at least
\[
 \max\!\left(0,
 \left\lceil\frac{q+1-2\sqrt q-B_p}{6}\right\rceil-m-s\right)
\]
terminal cubics.  A degree-\((m+3)\) locator has at most
\((m+3)!/6\) ordered marker--cubic decompositions.  Division by this
multiplicity gives the displayed leading term and error.
\end{proof}

For fixed \(r\), successive specialization of the selector gives a
deterministic construction using at most \((r-5)q\) partial-substitution
tests and \(q+1\) terminal pencil tests, assuming explicit coefficient access.
The dense selector has size exponential in \(r\); this is not a uniform
polynomial-time decoding theorem.

The next section combines this escape theorem with the explicit carrier and
the independent covering-radius inputs to classify the two highest coset-weight
shells.
